\section{Introduction}
\label{sec:intro}

Modern video diffusion transformers (DiTs) generate high-resolution video at
scale, but inference is typically instance-agnostic: the same frozen denoising
process is applied to every test clip, regardless of the motion, scene
layout, textures, or camera dynamics visible in the conditioning frames. In
video continuation, those conditioning frames are exactly the part of the
test input that carries instance-specific signal, yet that signal is not used
to specialize the model's denoising behavior at inference time.

Test-time adaptation (TTA) is the natural mechanism for exploiting this
signal, but the standard tools are poorly matched to a 13.6B-parameter video
DiT operating on a single short clip. Full-model tuning exposes billions of
parameters to a few seconds of supervision; LoRA-style adapters introduce
thousands to millions of free parameters that either overfit within a few
gradient steps or fail to move the model in a useful direction. What we want
is a TTA mechanism whose adaptation surface is small enough to be stable on a
single clip, yet expressive enough to perturb the model's behavior in a
direction the conditioning frames actually justify.

We introduce \method, an ultra-lightweight test-time adaptation method that
steers a frozen video DiT through its modulation pathway rather than through
weight updates. \method learns a compact additive residual $\delta$ in the
timestep embedding space and shares it across all transformer blocks. Each
block's frozen adaptive layer-normalization (adaLN) projection then maps the
shared residual into block-specific shift, scale, and gate vectors. The
trainable surface is global---one residual---but the induced perturbation is
per-block, because the frozen adaLN projections supply learned de-tying. This
\emph{structured weight tying} is what makes \method both compact (no new
layers, no source-data access at adaptation time) and expressive (every
block's modulation is perturbed independently). The residual is optimized
only on the observed conditioning frames using the standard denoising
objective, applied during generation, and discarded afterward.

\paragraph{Contributions.}
\begin{itemize}
  \item We formulate test-time adaptation for large video DiTs as
        timestep-pathway steering rather than direct weight adaptation.
  \item We introduce \method, a shared residual applied to the timestep
        embedding before per-block adaLN projections, and identify the
        structured weight-tying mechanism that lets a single trainable
        residual induce per-block modulation.
  \item We compare \method against No-TTA, LoRA at multiple ranks, and a
        TinyLoRA/SVD variant, on Panda-70M and UCF-101 with standard- and
        long-horizon evaluations at 1000-video scale.
  \item On standard-horizon Panda continuation, \method leaves FVD and
        per-frame metrics unchanged within noise relative to No-TTA at
        999 videos; on long-horizon Panda we
        report a controlled FVD/FID trade and treat long-horizon coherence
        as an open case.
  \item We document negative results---LoRA frequently trails baseline,
        aggressive adaptation overfits, several auxiliary regularizers do
        not improve the Pareto frontier---so that the modulation-pathway
        framing is not oversold.
\end{itemize}
